\documentclass{aastex631}

\usepackage{amsmath}
\usepackage{amssymb}
\usepackage{amsthm}
\usepackage{booktabs}
\usepackage{enumerate}

\newcommand{\blm}{b_{\ell m}}
\newcommand{\alm}{a_{\ell m}}
\newcommand{\Hlm}{H_{\ell m}}
\newcommand{\Dlm}{D^\ell_{mm'}}
\newcommand{\Tgnd}{T_\mathrm{gnd}}
\newcommand{\lmax}{\ell_\mathrm{max}}

\theoremstyle{plain}
\newtheorem{theorem}{Theorem}
\theoremstyle{remark}
\newtheorem{remark}{Remark}
\newtheorem{condition}{Condition}

% -------------------------------------------------------------------
\begin{document}

\title{Joint Fitting of Beam and Sky Using Rotations}

\correspondingauthor{Christian H. Bye}
\email{chbye@berkeley.edu}

\author{Christian H. Bye}
\author{Aaron R. Parsons}
\affiliation{Department of Astronomy, University of California Berkeley,
  Berkeley, CA 94720, USA}
\affiliation{Radio Astronomy Laboratory, University of California Berkeley,
  Berkeley, CA 94720, USA}

\begin{abstract}
We assess the theoretical basis for simultaneously recovering both the
antenna beam pattern $b_{\ell m}$ and the sky brightness temperature
$a_{\ell m}$ from a single rotating dipole antenna, without prior
knowledge of the beam.
The measurement equation is bilinear in $(b_{\ell m}, a_{\ell m})$,
so a single linear estimator cannot solve the stacked system.
Working in the spherical harmonic domain, we prove an identifiability
theorem: under full SO(3) rotation coverage, a known horizon profile,
and a known ground temperature $\Tgnd$, the beam and sky are uniquely
recoverable for all $\ell \geq 0$ via a constructive three-step
estimator --- Wigner projection, rank-1 SVD, and a linear system for
the per-$\ell$ amplitude scales.
The proof is constructive and doubles as the estimation algorithm.
We then analyze the four conditions of the theorem in detail as they
apply to the experiment.
The most important is incomplete SO(3) coverage: at geographic latitude
$L$ the three available rotation axes (mechanical tilt, turntable, and
Earth rotation) generate but do not uniformly sample SO(3), producing
cross-$\ell$ leakage whose magnitude and structure can be computed
analytically from the known rotation coverage.
We show that incomplete coverage calls for data-weighted Wigner
projections that optimally trade leakage against noise, and that
the weighting depends on LST in a non-trivial way.
The ground temperature condition is addressed by switched-load
absolute calibration (as implemented in EIGSEP), under the assumption
of spatially uniform ground emission.
Horizon profile errors propagate linearly into beam errors via a
known transfer matrix, providing a direct systematic budget.
We identify which horizon geometries guarantee a full-rank constraint
system, a question of independent theoretical interest.
\end{abstract}

% -------------------------------------------------------------------
\section{Experiment Setup} \label{sec:setup}

\subsection{Measurement Equation}

The baseline measurement equation for a single dipole is
\begin{equation} \label{eq:linear}
  \mathbf{y} = A\mathbf{x} + \mathbf{n},
\end{equation}
where $\mathbf{y}$ is the data vector (measured brightness temperature
as a function of time), $A$ is the design matrix encoding the antenna
beam pattern, pointing, and horizon profile, $\mathbf{x}$ is the sky
vector (brightness temperature as a function of pixel or SH
coefficient), and $\mathbf{n}$ is noise following the radiometer
equation.
With perfect knowledge of $A$, Equation~\ref{eq:linear} is solved
optimally by a Wiener filter.

\subsection{Mechanical Rotation Setup}

The antenna is mounted on a box with two mechanical rotation axes:
\begin{itemize}
  \item \textbf{X-axis (suspension/tilt axis):} tilts the antenna,
    sweeping the horizon cut across the beam pattern.
    Range $[0, \pi]$ sufficient for physical coverage.
  \item \textbf{Z$'$-axis (body-fixed turntable):} rotates the antenna
    in azimuth relative to the box. Range $[0, 2\pi]$.
\end{itemize}
Earth rotation provides a continuous third rotation over a 24-hour
sidereal day.
At geographic latitude $L$, the celestial pole in the local horizontal
frame sits at
\begin{equation}
  \hat{n}_\mathrm{pole} = \cos(L)\,\hat{n}_\mathrm{N} + \sin(L)\,\hat{z},
\end{equation}
where $\hat{z}$ is the local zenith and $\hat{n}_\mathrm{N}$ is
horizontal north; Earth rotation is around $\hat{n}_\mathrm{pole}$,
not around $\hat{z}$.
The three axes are summarized in Table~\ref{tab:axes}.

\begin{table}
  \centering
  \begin{tabular}{lll}
    \toprule
    \textbf{Rotation} & \textbf{Axis (local frame)} & \textbf{Range} \\
    \midrule
    X-tilt (mechanical)
      & $\hat{x}$ (horizontal)
      & $[0, \pi]$ \\
    Z$'$-turntable (mechanical)
      & $\hat{z}$ (zenith)
      & $[0, 2\pi]$ \\
    Earth/LST
      & $\cos(L)\,\hat{n}_\mathrm{N} + \sin(L)\,\hat{z}$
      & $[0, 2\pi]$ \\
    \bottomrule
  \end{tabular}
  \caption{The three rotation axes in the local horizontal frame.
    The Earth rotation axis coincides with the zenith only at the pole
    ($L=90^\circ$) and is purely horizontal at the equator ($L=0^\circ$),
    where it becomes parallel to the turntable and provides no
    independent coverage.
    Coverage quality is maximized at intermediate latitudes.}
  \label{tab:axes}
\end{table}

% -------------------------------------------------------------------
\section{The Problem is Bilinear} \label{sec:bilinear}

In the spherical harmonic (SH) domain, beam and sky expand as
\begin{equation}
  B(\hat{n}) = \sum_{\ell m} \blm\, Y_{\ell m}(\hat{n}), \qquad
  T(\hat{n}) = \sum_{\ell m} \alm\, Y_{\ell m}(\hat{n}).
\end{equation}
A rotation $R$ acts on the beam via Wigner $D$-matrices, giving the
full measurement
\begin{equation} \label{eq:SH}
  y(R) = \underbrace{\sum_\ell \sum_{m,m'} \blm\, \Dlm(R)\, a_{\ell m'}}_{\text{sky term, bilinear in }(b,a)}
        + \underbrace{\Tgnd \sum_{\ell m} \blm\, \Hlm(R)}_{\text{ground term, linear in }b}
        + n,
\end{equation}
where $\Dlm(R)$ is fully determined by the known rotation angle,
\begin{equation}
  \Hlm(R) = \int_{\hat{n}\,\in\,\mathrm{gnd}(R)} Y_{\ell m}(\hat{n})\, d\Omega
\end{equation}
is the integral of $Y_{\ell m}$ over the lower hemisphere at tilt $R$
(analytically computable from the horizon profile), and the ground term
is \emph{linear in $b$ only} because the horizon mask is entirely
determined by the known rotation and horizon profile.

Stacking observations at $N_R$ rotations gives
\begin{equation}
  \tilde{\mathbf{y}} = \tilde{A}(\mathbf{b})\,\mathbf{a} + \tilde{\mathbf{n}},
\end{equation}
which is \textbf{bilinear} in $(\mathbf{b}, \mathbf{a})$.
A single Wiener filter cannot solve it.
The question is whether, with sufficient rotation diversity and a known
horizon, the system is nevertheless \emph{identifiable} --- and if so,
what the explicit estimator is.

% -------------------------------------------------------------------
\section{Identifiability Theorem} \label{sec:theorem}

\subsection{Intuition: The Wigner Projection as an SO(3) Fourier Transform}

Before stating the theorem, we develop the key tool --- the Wigner
projection --- by analogy with the Fourier transform on a circle.

\textbf{Fourier analogy on SO(2).}
Suppose you observe $y(\phi) = f(\phi) \cdot g(\phi)$ on a circle,
where $f$ and $g$ are unknown functions of the rotation angle $\phi$.
A single measurement cannot separate them.
But if you rotate by many angles $\phi_i$ and form the weighted sum
\begin{equation}
  \sum_i w_i\, e^{-im\phi_i}\, y(\phi_i) \approx \hat{f}_m \cdot \hat{g}_0,
\end{equation}
the Fourier orthogonality $\sum_i w_i e^{i(m-m')\phi_i} = \delta_{mm'}$
projects out the single product $\hat{f}_m \cdot \hat{g}_0$.
By varying $m$, you recover all Fourier coefficients of $f$ (up to the
overall scale set by $\hat{g}_0$), and the full outer product
$\hat{f}_m \hat{g}_{m'}$ by varying the rotation axis.

\textbf{Lifting to SO(3).}
On the sphere, the rotation group is SO(3) and the role of $e^{im\phi}$
is played by the \emph{Wigner D-matrices} $D^\ell_{mm'}(R)$.
For a rotation $R$ with ZYZ Euler angles $(\alpha, \beta, \gamma)$:
\begin{equation}
  D^\ell_{mm'}(\alpha, \beta, \gamma)
  = e^{-im\alpha}\, d^\ell_{mm'}(\beta)\, e^{-im'\gamma},
\end{equation}
where $d^\ell_{mm'}(\beta)$ is a real matrix (the Wigner small-$d$
matrix) whose entries are polynomials in $\cos(\beta/2)$ and
$\sin(\beta/2)$.
For $\ell=1$ these are simply the standard rotation matrix elements.

The SO(3) analogue of Fourier orthogonality is the Wigner
orthogonality relation (Equation~\ref{eq:wigner_ortho}):
\begin{equation*}
  \sum_i w_i\, \overline{D^{\ell_0}_{m_0 m'_0}(R_i)}\,
  D^\ell_{mm'}(R_i)
  \approx \frac{8\pi^2}{2\ell+1}\,
  \delta_{\ell\ell_0}\,\delta_{mm_0}\,\delta_{m'm'_0}.
\end{equation*}
Multiplying the data $y(R_i)$ by $\overline{D^{\ell_0}_{m_0 m'_0}(R_i)}$
and summing over rotations therefore projects out the single product
$b_{\ell_0 m_0} \cdot a_{\ell_0 m'_0}$ from the bilinear sum.
This is the \textbf{Wigner projection}: the SO(3) analogue of taking
a Fourier coefficient.

\textbf{The resulting matrix.}
Performing the Wigner projection for all $(m_0, m'_0)$ at a fixed $\ell_0$
assembles the $(2\ell_0+1)\times(2\ell_0+1)$ matrix:
\begin{equation}
  [\hat{X}_{\ell_0}]_{m_0,\,m'_0}
  = \frac{2\ell_0+1}{8\pi^2}
    \sum_i w_i\,\overline{D^{\ell_0}_{m_0 m'_0}(R_i)}\,y(R_i)
  \approx b_{\ell_0 m_0}\, a_{\ell_0 m'_0}.
\end{equation}
The right-hand side is the outer product
$\mathbf{b}_{\ell_0}\mathbf{a}_{\ell_0}^T$ --- a rank-1 matrix whose
SVD immediately yields the beam and sky mode vectors.
The Wigner projection converts the bilinear estimation problem into a
rank-1 matrix recovery problem, one $\ell$-block at a time.

\subsection{Theorem Statement and Proof}

We state and prove identifiability constructively.
The proof yields the estimation algorithm directly.

\begin{theorem}[Joint beam--sky identifiability]
Consider the measurement model~\ref{eq:SH}.
Under the following five conditions:
\begin{enumerate}
  \item[\textbf{C1.}] \textbf{SO(3) coverage:} rotations $\{R_i\}$ with
    quadrature weights $\{w_i\}$ approximate the Haar integral on SO(3),
  \item[\textbf{C2.}] \textbf{Known $\Tgnd$:} the ground emission
    temperature is known absolutely,
  \item[\textbf{C3.}] \textbf{Known horizon:} the horizon profile is
    known, so $\Hlm(R)$ is computable,
  \item[\textbf{C4.}] \textbf{Full-rank $C$-matrix:} the operator
    $C_{\ell_0 m_0 m'_0;\,\ell m}$ defined in Equation~\ref{eq:Cmat}
    is full rank,
  \item[\textbf{C5.}] \textbf{Non-degenerate signals:}
    $b_{\ell m} \neq 0$ and $a_{\ell m'} \neq 0$ for modes of interest,
  \item[\textbf{C6.}] \textbf{Sky asymmetry:} the sky has nonzero
    $m' \neq 0$ power at each multipole order $\ell$ of interest,
    i.e., the sky is not perfectly axisymmetric at any $\ell$,
\end{enumerate}
the beam $\blm$ and sky $\alm$ are \textbf{uniquely identifiable}
for all $\ell \geq 0$.
\end{theorem}

\begin{proof}
The proof is constructive in three steps.
We first address a subtlety in Step 1 that the theorem statement
glosses over: the raw data $y(R_i)$ contains both the sky term and
the ground term, so the Wigner projection of $y$ does not directly
give $X_{\ell_0}$.
We describe the practical procedure for isolating the sky term.

\textbf{Step 1: Extracting $X_\ell = \mathbf{b}_\ell
\mathbf{a}_\ell^T$ from the raw data.}

Applying the Wigner projection to the full raw data $y(R_i)$ gives
\begin{equation} \label{eq:proj_raw}
  \hat{Y}_{\ell_0;\,m_0 m'_0}
  \equiv \frac{2\ell_0+1}{8\pi^2}
  \sum_i w_i\, \overline{D^{\ell_0}_{m_0 m'_0}(R_i)}\, y(R_i)
  = b_{\ell_0 m_0}\, a_{\ell_0 m'_0}
  + \Tgnd \sum_{\ell m} b_{\ell m}\, G^{\ell_0}_{m_0 m'_0;\,\ell m}
  + \tilde{n},
\end{equation}
where
\begin{equation}
  G^{\ell_0}_{m_0 m'_0;\,\ell m}
  \equiv \frac{2\ell_0+1}{8\pi^2}
  \sum_i w_i\,\overline{D^{\ell_0}_{m_0 m'_0}(R_i)}\,\Hlm(R_i)
\end{equation}
is the Wigner projection of the ground term --- a known matrix under
C1 and C3, but a contaminant in the recovery of $X_{\ell_0}$.
The ground term must be removed before the SVD of Step 1 is valid.
We exploit the fact that the ground term is constant in LST
(the horizon is fixed in the local frame) while the sky term varies.

\textit{Sub-step 1a: Recover $m'_0 \neq 0$ columns of $X_{\ell_0}$
from the LST residual.}
Define the LST-varying residual
\begin{equation} \label{eq:resid}
  \delta y(R_\mathrm{mech}, t)
  \equiv y(R_\mathrm{mech}, t) - \bar{y}(R_\mathrm{mech}),
\end{equation}
where the overbar denotes the mean over LST samples at fixed
mechanical rotation $R_\mathrm{mech}$.
The ground term, being constant in LST, cancels exactly:
\begin{equation}
  \delta y(R_\mathrm{mech}, t)
  = \delta f_\mathrm{sky}(R_\mathrm{mech}, t) + \delta n(t).
\end{equation}
The Wigner projection of $\delta y$ using the full rotation
$R_i = R_\mathrm{mech} \cdot R_\mathrm{LST}(t)$ gives
\begin{equation} \label{eq:substep1a}
  \frac{2\ell_0+1}{8\pi^2}
  \sum_i w_i\,\overline{D^{\ell_0}_{m_0 m'_0}(R_i)}\,\delta y(R_i)
  = b_{\ell_0 m_0}\, a_{\ell_0 m'_0}
  + \underbrace{\sum_{\ell \neq \ell_0}\sum_{mm'}
      b_{\ell m} a_{\ell m'}\,\epsilon^{(\ell_0\ell)}_{m_0 m'_0;\,mm'}
    }_{\text{(i) cross-}\ell\text{ leakage}}
  + \underbrace{\sum_{\ell m} b_{\ell m}\,\eta^{(\ell)}_{m m'_0}\,
      a_{\ell 0}
    }_{\text{(ii) LST sampling}}
  + \underbrace{\tilde{n}_{\ell_0;\,m_0 m'_0}
    }_{\text{(iii) noise}}
  \quad \text{for } m'_0 \neq 0,
\end{equation}
where the three error terms are:
\begin{enumerate}[(i)]
  \item \textbf{Cross-$\ell$ leakage:}
    $\epsilon^{(\ell_0\ell)}_{m_0 m'_0;\,mm'} \equiv
    \frac{2\ell_0+1}{8\pi^2}\sum_i w_i\,
    \overline{D^{\ell_0}_{m_0 m'_0}(R_i)}\,D^\ell_{mm'}(R_i)
    - \delta_{\ell\ell_0}\delta_{mm_0}\delta_{m'm'_0}$
    is the deviation of the discrete quadrature from Haar measure ---
    exactly the coverage kernel $K^{\ell_0\ell} - \delta_{\ell\ell_0}\mathbf{I}$.
    It is a \emph{known bias}, computable from the rotation grid
    alone before any data are taken, and is zero for an exact
    SO(3) $t$-design with $t \geq 2\lmax$.

  \item \textbf{LST sampling error:}
    $\eta^{(\ell)}_{mm'} \equiv
    \frac{1}{N_t}\sum_j D^\ell_{mm'}(R_\mathrm{LST}(t_j))
    - d^\ell_{m0}\,\delta_{m'0}$
    is the deviation of the discrete LST average from the
    complete sidereal average.
    For $N_t$ uniformly spaced LST samples it is
    $\mathcal{O}(1/N_t)$ for $\ell \lesssim N_t/2$ and grows for
    higher $\ell$.
    It re-introduces a small $m'=0$ sky contamination into
    $\delta y$ and is also a \emph{known bias} given the LST
    sampling schedule.

  \item \textbf{Noise:}
    $\tilde{n}_{\ell_0;\,m_0 m'_0}$ has zero mean and variance
    $\sigma_n^2 / N_\mathrm{obs}$ (for uniform weights and
    Haar-distributed rotations), where
    $\sigma_n^2 = T_\mathrm{sys}^2/(\Delta\nu\,\tau)$.
    It is the only \emph{random} error term; errors (i) and (ii)
    are deterministic biases.
\end{enumerate}
Both bias terms (i) and (ii) can be corrected using the known
rotation and LST grids, at the cost of some noise amplification.
Uncorrected, they bias the SVD of $\hat{X}_{\ell_0}$ and hence
the beam and sky shape estimates.
This recovers the $m'_0 \neq 0$ columns of $X_{\ell_0}$
without ground contamination.
The $m'_0 = 0$ column is absent because the LST average removes
the axisymmetric ($m'=0$) sky modes along with the ground term.

\textit{Sub-step 1b: Recover the $m'_0 = 0$ column iteratively.}
A preliminary beam shape estimate $\hat{\mathbf{b}}^{(0)}_{\ell_0}$
is obtained from the partial SVD of the $m'_0 \neq 0$ block of
$X_{\ell_0}$ recovered in Sub-step 1a.
The $m'_0 = 0$ column of $X_{\ell_0}$ is then recovered from the
raw-data projection~\ref{eq:proj_raw} by subtracting the known
ground contamination:
\begin{equation} \label{eq:m0col}
  \hat{X}_{\ell_0;\,m_0, 0}
  = \hat{Y}_{\ell_0;\,m_0, 0}
  - \Tgnd \sum_{\ell m} \hat{b}^{(0)}_{\ell m}\,
    G^{\ell_0}_{m_0, 0;\,\ell m}.
\end{equation}
With the full $X_{\ell_0}$ matrix now assembled, its SVD gives
refined estimates $\hat{\mathbf{b}}^{(1)}_{\ell_0}$ and
$\hat{\mathbf{a}}^{(1)}_{\ell_0}$.
Substituting $\hat{\mathbf{b}}^{(1)}$ back into
Equation~\ref{eq:m0col} and repeating gives a convergent iteration.

The convergence condition is \emph{not} that $\Tgnd \ll T_\mathrm{sky}$
--- which would fail at $\nu \gtrsim 150$\,MHz where the ground
dominates.
The actual requirement is that the beam shape
$\hat{\mathbf{b}}^{(0)}_{\ell_0}$, estimated from the $m'_0 \neq 0$
block in Sub-step 1a, is accurate enough to predict the ground
contamination in Equation~\ref{eq:m0col}.

The $m'_0 \neq 0$ block of $X_{\ell_0}$ is
\begin{equation}
  X_{\ell_0}[\,:\,,\, m' \neq 0]
  = \mathbf{b}_{\ell_0}\, [a_{\ell_0 m'}]_{m'\neq 0}^T,
\end{equation}
which is rank-1 with left singular vector $\mathbf{b}_{\ell_0}$
\emph{if and only if the sky has nonzero $m' \neq 0$ power at order
$\ell_0$}.
This is a condition on the sky (condition C6 below), not on the
count of entries: the rank-1 structure means the $m' \neq 0$ block
contains only $4\ell_0$ independent numbers regardless of its size,
and a single nonzero column already determines $\mathbf{b}_{\ell_0}$
up to the overall $\alpha_{\ell_0}$ scale.
Multiple nonzero columns improve noise resistance but do not add
identifiability.

The quality of $\hat{\mathbf{b}}^{(0)}_{\ell_0}$ --- and hence the
accuracy of the Sub-step 1b correction --- scales with the SNR of the
$m' \neq 0$ entries:
\begin{equation}
  \mathrm{SNR}_{\hat{b}_{\ell_0}}
  \sim
  \frac{\|\mathbf{b}_{\ell_0}\|\,\|[a_{\ell_0 m'}]_{m'\neq 0}\|}
       {\sigma_\mathrm{noise}}.
\end{equation}
This degrades at high $\ell_0$ (faint small-scale sky modes) and
at high frequency (faint sky overall), setting a practical frequency
ceiling on the method for a given noise level.

\begin{remark}
At high frequencies ($\nu \gtrsim 150$\,MHz), $\Tgnd \gg T_\mathrm{sky}$
and the LST residual $\delta y$ is small in absolute amplitude but
remains a clean, uncontaminated sky signal.
The ground term is large and nearly constant in LST, so it is
precisely measured --- favorable for Step 2.
The binding constraint at high frequency is the SNR of the $m' \neq 0$
sky modes in $\delta y$, which determines beam shape quality
in Sub-step 1a.
The optimal Wigner projection weights are frequency-dependent and
distinct for the sky and ground projections: sky weights should
up-weight LST samples with maximal Galactic plane contribution,
while ground weights are nearly uniform in LST.
Both sets are computable analytically from the known rotation grid
and sky/ground SNR at each frequency.
\end{remark}

\textit{Summary of Step 1 output.}
The result is the rank-1 matrix
\begin{equation}
  X_{\ell_0} = \mathbf{b}_{\ell_0}\,\mathbf{a}_{\ell_0}^T
  \;\in\; \mathbb{C}^{(2\ell_0+1)\times(2\ell_0+1)},
\end{equation}
whose SVD gives the \emph{shapes} $\hat{\mathbf{b}}_{\ell_0}$ and
$\hat{\mathbf{a}}_{\ell_0}$ up to an unknown positive scale
$\alpha_{\ell_0}$:
$\mathbf{b}_{\ell_0} = \alpha_{\ell_0}\hat{\mathbf{b}}_{\ell_0}$,
$\mathbf{a}_{\ell_0} = \alpha_{\ell_0}^{-1}\hat{\mathbf{a}}_{\ell_0}$.
There is one such unknown per $\ell_0$, exactly
$\mathcal{G}_\mathrm{sky} \cong (\mathbb{R}^+)^{\lmax}$.

\textbf{Step 2: Ground term linear system solves for $\alpha_\ell$.}
Apply the Wigner projection to the ground term alone (i.e., to
$\hat{Y}_{\ell_0}$ minus the sky projection recovered in Step 1):
\begin{equation}
  \frac{2\ell_0+1}{8\pi^2}
  \sum_i w_i\, \overline{D^{\ell_0}_{m_0 m'_0}(R_i)}\, g(R_i)
  = \Tgnd \sum_{\ell m} b_{\ell m}\, C_{\ell_0 m_0 m'_0;\,\ell m},
\end{equation}
where
\begin{equation} \label{eq:Cmat}
  C_{\ell_0 m_0 m'_0;\,\ell m}
  \equiv \frac{2\ell_0+1}{8\pi^2}
  \sum_i w_i\,
  \overline{D^{\ell_0}_{m_0 m'_0}(R_i)}\,\Hlm(R_i)
\end{equation}
is a \emph{known matrix} computable from the horizon geometry and
rotation coverage alone (conditions C1, C3).
Substituting $b_{\ell m} = \alpha_\ell\,\hat{b}_{\ell m}$ from Step 1
gives a \emph{linear system in $\{\alpha_\ell\}$}:
\begin{equation} \label{eq:alphasystem}
  \sum_\ell \alpha_\ell \underbrace{\left(
    \Tgnd \sum_m \hat{b}_{\ell m}\,
    C_{\ell_0 m_0 m'_0;\,\ell m}
  \right)}_{M_{\ell_0 m_0 m'_0;\,\ell}}
  = \text{measured ground projection}.
\end{equation}
Under condition C4, $M$ is full rank and $\boldsymbol{\alpha}$ is
uniquely determined.

\textbf{Step 3: Monopole recovery.}
At $\ell_0 = 0$, $D^0_{00}(R) = 1$ for all $R$, so the Wigner
projection is simply the mean over all observations.
Sub-step 1a gives $b_{00}\,a_{00}$ from the mean sky term.
The ground projection at $\ell_0 = 0$ gives
$\Tgnd\,b_{00}\,H_{00} + \Tgnd\sum_{\ell>0} b_{\ell 0}\,H_{\ell 0}$,
where the second sum is known from Step 2.
This pins $b_{00}$ (since $H_{00}$ and $\Tgnd$ are known under C2,
C3), and the product $b_{00}\,a_{00}$ from Sub-step 1a then gives
$a_{00}$.
\end{proof}

\begin{remark}
This is a stronger result than it might appear.
The sky monopole $a_{00}$ --- the global mean brightness temperature,
which carries the 21-cm cosmological signal --- \emph{is} recoverable,
without an absolute load calibrator, \emph{provided} $\Tgnd$
is known absolutely (C2) and the horizon is known (C3).
An absolute calibrator is needed only to the extent that C2 is
imperfectly satisfied.
\end{remark}

% -------------------------------------------------------------------
\section{Analysis of the Theorem Conditions} \label{sec:conditions}

Each condition C1--C5 can fail in practice.
We analyze each in turn, with emphasis on C1 (SO(3) coverage) as the
most important for experimental design, and C2--C3 as the primary
systematic uncertainties.

% -------------------------------------------------------------------
\subsection{C1: SO(3) Coverage and the Cost of Incompleteness}
\label{sec:so3}

This is the most important condition for experimental design and
deserves detailed treatment.

\subsubsection{Coverage structure at latitude $L$}

The three rotation axes (Table~\ref{tab:axes}) span $\mathbb{R}^3$
for $0 < L < 90^\circ$ and any mount orientation not aligned
north--south, so they \emph{generate} all of SO(3).
However, generation is not the same as uniform sampling.
At the equator ($L=0^\circ$) the polar axis is horizontal and parallel
to the turntable, leaving only two independent axes.
At the pole ($L=90^\circ$) the polar axis coincides with the zenith
and again collapses to two axes.
Coverage is best at intermediate latitudes; at Berkeley
($L\approx38^\circ$), $\cos L \approx 0.79$ and Earth rotation sweeps
a large, independent arc.

\subsubsection{Leakage from incomplete coverage}

Under non-Haar sampling, the Wigner projection~\ref{eq:proj} acquires
cross-$\ell$ leakage.
Define the \emph{coverage kernel}
\begin{equation}
  K^{\ell_0 \ell}_{\,m_0 m'_0;\,m m'}
  \equiv \frac{2\ell_0+1}{8\pi^2}
  \sum_i w_i\,
  \overline{D^{\ell_0}_{m_0 m'_0}(R_i)}\,
  D^\ell_{mm'}(R_i),
\end{equation}
so that the projected measurement is
\begin{equation} \label{eq:leaky}
  \hat{X}_{\ell_0;\,m_0 m'_0}
  = \sum_{\ell mm'} K^{\ell_0\ell}_{\,m_0 m'_0;\,mm'}\,
    b_{\ell m}\, a_{\ell m'}.
\end{equation}
Under Haar measure, $K^{\ell_0\ell} = \delta_{\ell\ell_0}$
(no leakage).
For a finite, non-uniform grid, $K^{\ell_0\ell}$ is a known matrix
that can be computed analytically from the rotation grid
$\{R_i, w_i\}$.
The leakage terms ($\ell \neq \ell_0$) in Equation~\ref{eq:leaky}
act as a known additive systematic in each projected measurement.

\subsubsection{Optimal data weighting and LST dependence}
\label{sec:weighting}

Since the coverage kernel $K$ is fully computable, the leakage
is not a fundamental limitation --- it is a \emph{known bias} that
can be corrected.
The natural approach is to choose weights $\{w_i\}$ to simultaneously
minimize leakage and noise amplification.
Crucially, the sky and ground Wigner projections require
\emph{different} optimal weights, and both sets are
frequency-dependent.

\textbf{Sky projection weights.}
The sky term is accessed via the LST residual $\delta y(R,t)$
(Sub-step 1a of the proof), whose SNR is set by
$T_\mathrm{sky}/T_\mathrm{noise}$ — a ratio that falls with
frequency as $\nu^{-2.5}$.
The sky weights should up-weight LST samples where the Galactic
plane contributes maximally to $\delta y$ (high-$\ell$ sky modes)
and down-weight featureless sky samples.
The optimization is
\begin{equation} \label{eq:optweights_sky}
  \min_{\{w^{\rm sky}_i\}} \left[
    \lambda \sum_{\ell \neq \ell_0} \|K^{\ell_0\ell}\|_F^2
    + (1-\lambda) \sum_i \frac{(w^{\rm sky}_i)^2}{\mathrm{SNR}^2_{\delta y,\,i}}
  \right]
\end{equation}
subject to $K^{\ell_0\ell_0} = \mathbf{I}$.

\textbf{Ground projection weights.}
The ground term is accessed via the LST mean $\bar{y}(R)$
(or equivalently via $\hat{Y} - X$ after Step 1), whose SNR
is set by $T_\mathrm{gnd}/T_\mathrm{noise}$ — spectrally flat and
large at all frequencies.
At high frequency ($\nu \gtrsim 150$\,MHz), $T_\mathrm{gnd}$ dominates
$y(R)$ and the ground projection SNR is very high; the ground weights
can be nearly uniform.
The same optimization applies with $\mathrm{SNR}^2_{\delta y}$
replaced by $\mathrm{SNR}^2_{\bar{y}}$.

\textbf{Joint weighting.}
In practice, a single set of weights is used for the full observation
$y(R_i)$, but the sky and ground projections are formed from
different linear combinations of those observations ($\delta y$ and
$\bar{y}$ respectively) and so effectively use different weights.
The key point is that the LST structure matters:
at low frequency the Galactic plane transit is the dominant
feature and LST sampling near the transit should be dense;
at high frequency the ground term is the dominant signal and
LST uniformity matters more for the sky term but less overall.

\begin{remark}
At the limit of complete SO(3) coverage, both weight sets converge
to the Haar measure and the distinction vanishes --- the sky and
ground terms decouple exactly and both are recovered simultaneously
at any single LST instant.
With incomplete coverage, the two weight sets diverge, and
optimizing them independently squeezes the most information
from the available rotation grid at each frequency.
\end{remark}

% -------------------------------------------------------------------
\subsection{C2: Known Ground Temperature}
\label{sec:tgnd}

The identifiability theorem requires $\Tgnd$ to be known absolutely,
since it sets the scale of the linear ground-term constraints.
An error $\varepsilon$ in $\Tgnd$ propagates into an error in
$\boldsymbol{\alpha}$ via the linear system~\ref{eq:alphasystem}, and
thence into errors in both beam and sky.

EIGSEP implements switched-load absolute calibration to address this, using a
switched noise source and ambient temperature load to establish an
absolute receiver temperature scale.
For this calibration to pin $\Tgnd$, one key assumption is required:

\textbf{Spatial uniformity of ground emission.}
The calibration measures the total ground contribution at a known
effective temperature, but the ground term in Equation~\ref{eq:SH}
is a \emph{beam-weighted} integral over the visible ground.
If the ground temperature varies spatially, the effective $\Tgnd$ is
beam-shape-dependent and rotation-dependent, and a single scalar
$\Tgnd$ is an approximation.
The assumption required is that the ground is spatially uniform in
brightness temperature, at least over the footprint of the beam
sidelobes.
Violations --- dry vs.\ wet soil patches, vegetation, water bodies,
near-field scattering objects --- will produce an effective
$\Tgnd(R)$ that varies with rotation angle, introducing a systematic
not captured by the scalar calibration.

The severity of this systematic can in principle be bounded by
measuring $y$ at the zenith-pointing position (tilt $= 0$, all ground
below horizon) across all turntable azimuths: any azimuthal variation
in the ground-only measurement directly reflects spatial
non-uniformity of the ground.
This is a concrete diagnostic available from the data.

% -------------------------------------------------------------------
\subsection{C3: Known Horizon Profile and Error Propagation}
\label{sec:horizon}

\subsubsection{Role of the horizon}

With a known horizon profile, each observation at rotation $r$ and
local sidereal time $t$ splits as
\begin{equation} \label{eq:split}
  y_r(t) =
  \underbrace{\sum_{\ell mm'} \blm\,\Dlm(R)\,a_{\ell m'}}_{\text{bilinear}}
  +
  \underbrace{\Tgnd \sum_{\ell m} \blm\,\Hlm(R)}_{\text{linear in }b}
  + n_r(t).
\end{equation}
The ground term is purely linear in $b$ because the horizon mask is
known.
For a fixed mechanical rotation, the ground term is constant in LST
while the sky term varies, allowing regression to separate the two
contributions --- though for a wide beam, both terms change
substantially with tilt, and they must be fitted jointly (see
Section~\ref{sec:widbeam}).

\subsubsection{Radon-transform interpretation and wide-beam failure}
\label{sec:widbeam}

The naive knife-edge / Radon-transform strategy --- differentiate $y$
with respect to tilt $\phi$ to isolate a line integral of $b$ along
the horizon, then invert --- fails for a wide beam.
The derivative $\partial y/\partial \phi$ receives two contributions:
\begin{align}
  \frac{\partial y}{\partial \phi} &=
  \Tgnd
  \underbrace{\sum_{p\,\in\,\partial\mathrm{horizon}} b_p}_{\text{Radon term}}
  +
  \underbrace{\frac{\partial}{\partial \phi}
    \sum_{p\,\in\,\mathrm{sky}} b_p\,
    x(\alpha_p+\phi+\delta t)}_{\text{sky contamination (comparable for wide beam)}}.
\end{align}
The sky contamination cannot be removed without already knowing the
sky.
The constructive estimator of Section~\ref{sec:theorem} handles this
correctly by treating the full bilinear system.

\subsubsection{Horizon error propagation}

Suppose the true horizon profile differs from the assumed profile by
$\delta H(\phi)$, inducing an error $\delta \Hlm(R)$ in the computed
hemisphere integrals.
This perturbs the $C$-matrix in Equation~\ref{eq:Cmat} by
\begin{equation}
  \delta C_{\ell_0 m_0 m'_0;\,\ell m}
  = \frac{2\ell_0+1}{8\pi^2}
  \sum_i w_i\,
  \overline{D^{\ell_0}_{m_0 m'_0}(R_i)}\,\delta\Hlm(R_i),
\end{equation}
which is a \emph{known linear map} from horizon errors to $C$-matrix
errors.
The resulting error in $\boldsymbol{\alpha}$ follows from perturbing
the linear system~\ref{eq:alphasystem}:
\begin{equation} \label{eq:alphaerror}
  \delta\boldsymbol{\alpha} = -M^{-1}\,\delta M\,\boldsymbol{\alpha},
\end{equation}
where $\delta M$ is the perturbation to $M$ induced by $\delta C$.
This transfer matrix $M^{-1}\,\delta M$ is computable from the
nominal solution and provides a direct \emph{systematic error budget}:
the contribution of a horizon profile error of magnitude
$\|\delta H\|$ to the beam and sky recovery errors is bounded by
$\|M^{-1}\| \cdot \|\delta M\| \cdot \|\boldsymbol{\alpha}\|$.
A concrete numerical evaluation of this budget, using LiDAR-measured
horizon profiles for the EIGSEP site, is deferred to future work.

% -------------------------------------------------------------------
\subsection{C4: Full-Rank $C$-Matrix}
\label{sec:rank}

The $C$-matrix (Equation~\ref{eq:Cmat}) must be full rank for the
linear system~\ref{eq:alphasystem} to uniquely determine
$\boldsymbol{\alpha}$.
Since $\Hlm(R)$ mixes different $\ell$ (a tilted great-circle cut
through a hemisphere has power at all multipole orders), $C$ is
generically full rank.
However, special horizon geometries can cause rank deficiency.

Specifically, the $C$-matrix will be rank-deficient if and only if
there exists a non-zero vector $\boldsymbol{\alpha}$ such that
$\sum_\ell \alpha_\ell \Hlm(R) = 0$ for all $R$ and all
$(\ell_0, m_0, m'_0)$.
This requires the function $\sum_\ell \alpha_\ell \Hlm(R)$ to vanish
identically under Wigner projection, which is the condition that
the $\Hlm$ at different $\ell$ are linearly dependent under the
rotation group action.

Three cases are worth noting:

\begin{enumerate}
  \item \textbf{Flat horizon (no tilt variation):} if the tilt range
    is zero (antenna always horizontal), $\Hlm(R)$ only samples
    hemisphere integrals at one elevation cut.
    The matrix $C$ has rank 1 in the $\ell$ direction and
    $\boldsymbol{\alpha}$ is undetermined.
    Some tilt range is essential.

  \item \textbf{Rotationally symmetric horizon:} if the horizon is
    flat and featureless (same elevation at all azimuths), then
    $\Hlm(R)$ is non-zero only for $m=0$ modes, and $C$ loses rank
    in the $m$ direction.
    Azimuthal structure in the horizon --- trees, buildings, terrain
    --- provides $m \neq 0$ power in $\Hlm$ and improves rank.

  \item \textbf{Generic non-flat horizon:} provided the tilt range
    spans more than a few degrees and the horizon has some azimuthal
    variation (which is true for any real site), $C$ is generically
    full rank.
    Characterizing the minimum tilt range required for full rank at
    a given $\lmax$, and relating it to the horizon power spectrum,
    would be of theoretical interest and is left for future work.
\end{enumerate}

% -------------------------------------------------------------------
\subsection{C5: Non-degenerate Signals}

This condition is benign in practice.
For $b_{\ell m} = 0$ at some $\ell$, the beam has no power at that
multipole and that scale is simply not measurable (nor needed).
For $a_{\ell m'} = 0$, the sky has no power at that multipole and
no beam mode at $\ell$ can be constrained by the sky term; the
ground term still constrains the beam.
The only genuine failure is $b_{00} = 0$ (zero-gain antenna) which
is unphysical.

% -------------------------------------------------------------------
\subsection{C6: Sky Asymmetry}
\label{sec:skyasym}

This condition is the most practically constraining at high $\ell$
and high frequency.
The beam shape $\hat{\mathbf{b}}_\ell$ is estimated from the
$m' \neq 0$ columns of $X_\ell$ (Sub-step 1a of the proof), which
requires the sky to have nonzero $m' \neq 0$ power at order $\ell$.
The condition is necessary: if the sky is perfectly axisymmetric at
some $\ell$ (all power in $m'=0$), the $m' \neq 0$ block vanishes
and the beam shape cannot be recovered from the LST residual at that
$\ell$.

For the real diffuse radio sky, C6 is satisfied for all $\ell$ up
to a practical $\lmax^\mathrm{sky}$, with the Galactic plane
providing strong $m' \neq 0$ power.
The quality of the beam shape estimate scales with the $m' \neq 0$
sky power at order $\ell$, which can be computed directly from the
GSM covariance matrix as a function of $\ell$ and frequency:
\begin{equation}
  \mathrm{SNR}_{\hat{b}_\ell}(\nu)
  \sim
  \frac{\|\mathbf{b}_\ell\|\,\|[a_{\ell m'}(\nu)]_{m'\neq 0}\|}
       {T_\mathrm{sys}/\sqrt{\Delta\nu\,\tau}},
\end{equation}
where the numerator is determined by the sky model and the
denominator by the radiometer noise.
This SNR can be evaluated as a function of $(\ell, \nu)$ before
taking data, providing a map of which beam modes are reliably
recoverable across the frequency band.

% -------------------------------------------------------------------
\section{Hierarchical Structure from Full LST Coverage}
\label{sec:hierarchy}

Time-averaging over a complete sidereal day yields a useful
decomposition that does not lose information but reorganizes it
for efficient estimation.
By Wigner orthogonality of the Earth-rotation component,
\begin{equation}
  \bigl\langle D^\ell_{mm'}(R_\mathrm{mech}\cdot R_\mathrm{LST}(t))
  \bigr\rangle_t = d^\ell_{m0}(\theta_R)\,\delta_{m'0},
\end{equation}
so the time-averaged observation is
\begin{equation} \label{eq:ybar}
  \bar{y}(R) = \sum_{\ell m} \blm\, d^\ell_{m0}(\theta_R)\, a_{\ell 0}
             + \Tgnd\sum_{\ell m}\blm\,\Hlm(R).
\end{equation}
Only the $m'=0$ sky modes $a_{\ell 0}$ appear in $\bar{y}$.
The sky collapses from $(L_s+1)^2$ unknowns to $(L_s+1)$ --- a
factor of $(L_s+1)$ reduction.
The LST-varying residual $\delta y = y - \bar{y}$ carries all
$m'\neq 0$ sky information but loses the ground term anchor.

This decomposition is lossless (it is a bijection on the data) and
is valid under complete 24-hr LST coverage.
With partial coverage, $\langle D^\ell_{mm'}\rangle_t \neq
d^\ell_{m0}\delta_{m'0}$ and cross-$m'$ leakage occurs, exactly
analogous to the cross-$\ell$ leakage from incomplete SO(3) coverage.
The optimal weighting of Section~\ref{sec:weighting} addresses both
simultaneously.

% -------------------------------------------------------------------
\section{Degrees of Freedom} \label{sec:dof}

The unknowns per frequency channel are:
beam $\blm$ for $\ell = 0,\ldots,L_b$: $(L_b+1)^2$ real DOF;
sky $\alm$ for $\ell = 0,\ldots,L_s$: $(L_s+1)^2$ real DOF.
The available observations are $N_R \times N_t$ real measurements.
The formal identifiability condition $N_R N_t \geq (L_b+1)^2 +
(L_s+1)^2$ is necessary but not sufficient; the identifiability
theorem provides the sufficient conditions.

Within each $\ell$-block, the sky term is a linear functional of the
rank-1 matrix $X_\ell = \mathbf{b}_\ell \mathbf{a}_\ell^T$.
By compressed sensing theory (RIP for rank-1 matrices under
Haar-random Wigner measurement operators), the number of measurements
needed to recover $X_\ell$ at order $\ell$ is
\begin{equation}
  N_\mathrm{meas}^\ell \;\gtrsim\; 4(2\ell+1) - 4 \approx 8\ell,
\end{equation}
setting a minimum rotation grid density per $\ell$-block.

% -------------------------------------------------------------------
\section{Estimation Strategy} \label{sec:estimation}

The constructive proof of the theorem directly yields the estimation
algorithm.

\textbf{Step 1 (Spectral initialization):}
Form the Wigner-projected matrices $\hat{X}_\ell$ per $\ell$-block
using the optimal weights from Equation~\ref{eq:optweights}.
Take the top rank-1 SVD of each $\hat{X}_\ell$ to get initial
estimates $\hat{\mathbf{b}}^{(0)}_\ell$, $\hat{\mathbf{a}}^{(0)}_\ell$.
This step is non-iterative and provides a globally consistent
initialization.

\textbf{Step 2 (Scale fixing):}
Solve the linear system~\ref{eq:alphasystem} for
$\boldsymbol{\alpha}$ using the ground-term Wigner projections.
This uniquely pins the per-$\ell$ amplitudes and completes the
closed-form estimator under ideal conditions.

\textbf{Step 3 (Joint refinement via ALS or Gibbs sampling):}
Run alternating least squares (ALS) or Gibbs sampling over all
$(t, R)$ simultaneously, initialized from $(\hat{\mathbf{b}}^{(0)},
\hat{\mathbf{a}}^{(0)})$.
Each ALS step fixes one block and applies a Wiener filter to the
other; the ground term enters as a known linear contribution.
For Gaussian noise and Gaussian priors, each Gibbs conditional is
Gaussian (a linear solve), and the sampler correctly propagates
uncertainty through the bilinear structure.

ALS and Gibbs both converge at linear rate $\rho_\mathrm{ALS} =
\|\Sigma_{bb}^{-1}\Sigma_{ba}\Sigma_{aa}^{-1}\Sigma_{ab}\|_2^{1/2}$,
the largest canonical correlation between beam and sky in the joint
Fisher matrix.
Convergence requires $\rho < 1$, which holds when the sky is
structured (Galactic plane), the beam is asymmetric, and informative
priors regularize the degenerate directions.
Counterintuitively, high SNR alone degrades convergence:
at high SNR the posterior concentrates on the bilinear manifold and
$\rho \to 1$; priors are essential.

% -------------------------------------------------------------------
\section{Relationship to Rotation-Free Experiments and the Prior Spectrum}
\label{sec:edges}

\subsection{What Rotation-Free Experiments Know About Their Beam}

Single-antenna experiments without rotation diversity measure
$T_\mathrm{ant}(\nu)$, the beam-weighted sky brightness temperature
as a function of frequency, at fixed pointing.
The beam enters as a nuisance: its chromaticity couples spatial sky
structure into the frequency spectrum, producing features that
mimic or obscure the 21-cm signal.
Such experiments handle this by computing the beam from an
electromagnetic simulation (FMCW or HFSS) and treating the result
as truth.
The beam is never measured in situ; the residual uncertainty between
the simulated and true beam is an unquantified systematic.

EIGSEP has the same HFSS simulation capability, and arguably a
better one: the canyon suspension places the antenna in a nearly
free-space electromagnetic environment, eliminating the ground
coupling to the antenna structure that complicates terrestrial
simulations.
EIGSEP therefore starts with the same level of beam knowledge as
rotation-free experiments.
The question rotis addresses is how much further the rotation data
can push that knowledge.

\subsection{The Prior Spectrum}

The rotis estimation pipeline operates along a continuous spectrum
of beam prior strength, parameterized by the fractional uncertainty
$\sigma_\ell$ in the beam coefficients $\blm$:

\begin{equation}
  \blm \sim \mathcal{N}\!\left(b^\mathrm{HFSS}_{\ell m},\;
  \sigma_\ell^2\right),
\end{equation}

where $b^\mathrm{HFSS}_{\ell m}$ is the HFSS-simulated beam.
The prior width $\sigma_\ell$ can vary with $\ell$, reflecting
different levels of simulation accuracy at different angular scales.

\begin{table}
  \centering
  \begin{tabular}{c|c|c|c}
    \toprule
    \textbf{Regime} & \textbf{$\sigma_\ell$} & \textbf{Beam free parameters}
      & \textbf{Equivalent to} \\
    \midrule
    Rotation-free
      & $\to 0$
      & 0 (simulation is truth)
      & Standard single-antenna pipeline \\
    \hline
    Perturbative
      & $1$--$5\%$
      & Low-$\ell$ corrections only
      & Single-antenna $+$ in-situ verification \\
    \hline
    Weakly informative
      & $\sim 50\%$
      & All modes, constrained
      & Rotation-assisted calibration \\
    \hline
    Free beam
      & $\to \infty$
      & $(L_b+1)^2$ free parameters
      & Full rotis theorem \\
    \bottomrule
  \end{tabular}
  \caption{The prior spectrum for beam estimation in rotis.
    The rotation-free limit ($\sigma_\ell \to 0$) corresponds to
    treating the EM simulation as truth without in-situ verification.
    The rotis identifiability theorem guarantees recovery in the
    $\sigma_\ell \to \infty$ limit under conditions C1--C6.
    The perturbative regime is the most immediately practical:
    it uses the simulation as initialization and measures corrections
    from the rotation data.}
  \label{tab:prior_spectrum}
\end{table}

In the perturbative regime, the effective number of beam free
parameters is small --- perhaps a handful of low-$\ell$ modes
representing gain uncertainty, beam ellipticity errors, and horizon
coupling residuals.
The SNR required to constrain these few parameters is easily
achievable with the EIGSEP noise budget, unlike the full free-beam
case which requires SNR $\sim 10^5$ per mode.

\subsection{The Key Scientific Output: Posterior Width vs.\ Prior Width}

The most important diagnostic of the rotis pipeline is the function
\begin{equation}
  \sigma^\mathrm{post}_\ell(\sigma^\mathrm{prior}_\ell),
\end{equation}
the posterior width on $\blm$ as a function of the prior width,
computed from the Gibbs sampler at each $\ell$.
This function directly answers: \emph{how much does the rotation
data constrain the beam beyond the HFSS simulation?}

Three qualitatively distinct behaviors are possible:
\begin{enumerate}
  \item $\sigma^\mathrm{post}_\ell \ll \sigma^\mathrm{prior}_\ell$:
    the data are strongly informative at this $\ell$; rotation
    diversity substantially improves on the simulation.
  \item $\sigma^\mathrm{post}_\ell \sim \sigma^\mathrm{prior}_\ell$:
    the data add nothing; the posterior tracks the prior.
    Either this $\ell$-mode is not constrained by the rotation
    grid, or the sky has insufficient $m' \neq 0$ power (C6).
  \item $\sigma^\mathrm{post}_\ell \sim \sigma^\mathrm{noise}_\ell$:
    the noise floor is reached; rotation data constrain this mode
    as well as the data allow regardless of prior.
\end{enumerate}
The crossover between regimes (1) and (2) as a function of $\ell$
defines the effective angular resolution of the in-situ beam
measurement.
The crossover between regimes (1) and (3) as a function of
$\sigma^\mathrm{prior}$ defines how informative the rotation data
are relative to the simulation.

\subsection{Hierarchy of Scientific Claims}

The prior spectrum framing supports a natural hierarchy of
increasingly ambitious claims, all derivable from the same dataset:

\begin{enumerate}
  \item \textbf{Beam verification (immediately achievable):}
    Use rotation data to test whether the HFSS simulation matches
    the in-situ beam response as a function of rotation angle.
    This is a null test unavailable to rotation-free experiments.
    Significant regardless of the 21-cm signal.

  \item \textbf{Perturbative beam correction (achievable with
    tight priors):}
    Use rotation data to measure and correct low-order beam
    perturbations around the HFSS simulation.
    Feed the corrected beam into a standard single-antenna spectral
    analysis.
    Reduces beam systematic uncertainty below the rotation-free level.

  \item \textbf{Rotation-assisted beam recovery (weakly informative
    prior):}
    Recover the beam with significant freedom from the simulation,
    using rotation diversity as the primary constraint.
    Demonstrates the in-principle power of the rotis approach.

  \item \textbf{Free-beam recovery (full theorem):}
    Recover beam and sky jointly with no simulation prior.
    Requires SNR $\sim 10^5$ per beam mode; aspirational with
    current noise budgets.
\end{enumerate}

Claims 1 and 2 are immediately achievable and publishable.
Claim 3 depends on the posterior-width diagnostic showing
informative constraints beyond the prior.
Claim 4 is the theoretical limit guaranteed by the identifiability
theorem.

% -------------------------------------------------------------------
\section{Fisher Information Summary} \label{sec:fisher}

A 1D elevation-slice numerical Fisher analysis (22 beam pixels over
$\pm 80^\circ$, 16 sky pixels over $1$--$75^\circ$, $\Tgnd = 300$~K)
gives the following qualitative results:

\begin{itemize}
  \item Tilt range is the primary driver of beam mode recovery.
    The eigenvalue spectrum of the joint FIM shows a ``knee'' that
    lifts with tilt range.
  \item LST diversity independently constrains sky modes.
    Without sidereal time variation, sky modes are severely
    underdetermined even with good beam constraints.
  \item Two to four modes remain near zero regardless of tilt and LST,
    corresponding to the conditions in Table~\ref{tab:conditions}.
  \item Structured sky (e.g., the Galactic plane) dramatically improves
    sky mode recovery.
  \item Beam pixels near the horizon edge are best constrained;
    pixels far from the horizon are constrained primarily via LST
    drift against sky structure.
\end{itemize}

% -------------------------------------------------------------------
\section{Summary} \label{sec:summary}

\begin{table}
  \centering
  \begin{tabular}{c|c|c|c}
    \toprule
    \textbf{Condition} & \textbf{What fails if violated}
      & \textbf{Mitigation} & \textbf{Status} \\
    \midrule
    C1: SO(3) coverage
      & Cross-$\ell$ leakage; biased $X_\ell$
      & Optimal data weighting; coverage kernel from known grid
      & Primary design driver \\[4pt]
    \hline
    C2: Known $\Tgnd$
      & $\boldsymbol{\alpha}$ biased; monopole error
      & Switched-load absolute calibration (EIGSEP);
        spatial uniformity assumption
      & Addressed by EIGSEP \\[4pt]
    \hline
    C3: Known horizon
      & $C$-matrix errors; beam/sky bias
      & LiDAR survey; transfer matrix $M^{-1}\delta M$
      & Budget TBD \\[4pt]
    \hline
    C4: Full-rank $C$
      & $\boldsymbol{\alpha}$ underdetermined
      & Non-zero tilt; azimuthally varied horizon
      & Generically satisfied \\[4pt]
    \hline
    C5: Non-zero signals
      & $\ell$-modes unrecoverable
      & N/A for physical antennas
      & Not a concern \\[4pt]
    \hline
    C6: Sky asymmetry
      & Beam shape unrecoverable from $\delta y$
      & GSM confirms $m'\neq 0$ power; SNR map in $(\ell,\nu)$
      & Frequency/\,$\ell$-dependent \\
    \bottomrule
  \end{tabular}
  \caption{Summary of the six identifiability conditions, the
    consequence of each failure, the mitigation strategy, and the
    current experimental status.}
  \label{tab:conditions}
\end{table}

The central results are:

\begin{enumerate}

  \item \textbf{Beam and sky are jointly identifiable}, including the
    sky monopole, under the six conditions of the theorem.
    The proof is constructive: Wigner projection recovers $X_\ell$,
    the ground term linear system recovers per-$\ell$ scales, and
    the monopole follows from the $\ell=0$ ground projection.

  \item \textbf{Incomplete SO(3) coverage introduces known,
    correctable bias.}
    The coverage kernel $K^{\ell_0\ell}$ is computable from the
    known rotation grid and produces a known leakage systematic.
    Optimal data weighting minimizes the leakage--noise tradeoff
    and depends on LST in a non-trivial way.

  \item \textbf{$\Tgnd$ uncertainty is the primary condition
    requiring external calibration.}
    Spatial uniformity of ground emission is the key assumption;
    its violation is diagnosable from azimuthal variation of the
    ground-only signal.

  \item \textbf{Horizon profile errors propagate linearly into
    beam and sky errors} via the known transfer matrix
    $M^{-1}\delta M$. A quantitative systematic budget requires
    a site-specific LiDAR survey.

  \item \textbf{Generic non-flat horizons produce full-rank $C$.}
    The minimum tilt range for full rank as a function of $\lmax$
    and horizon power spectrum is an open theoretical question.

  \item \textbf{The HFSS simulation provides the same beam
    knowledge as rotation-free experiments as a starting point.}
    The rotis pipeline operates along a prior spectrum from
    tight (rotation-free limit) to free (full theorem), parameterized by
    the fractional beam uncertainty $\sigma_\ell$.
    The key scientific output is the posterior width
    $\sigma^\mathrm{post}_\ell$ as a function of prior width
    $\sigma^\mathrm{prior}_\ell$, which directly measures how
    much the rotation data improve on the simulation at each
    angular scale (Table~\ref{tab:prior_spectrum}).
    Claims 1 and 2 of the scientific hierarchy
    (Section~\ref{sec:edges}) are achievable with current noise
    budgets; Claims 3 and 4 are aspirational.

\end{enumerate}

\subsection{Immediate Next Steps}

\begin{enumerate}
  \item Compute the coverage kernel $K^{\ell_0\ell}$ numerically
    for the planned rotation grid and quantify cross-$\ell$ leakage
    vs.\ $\ell_0$.
  \item Solve the optimal weighting problem
    (Equation~\ref{eq:optweights_sky}) for the coverage and assess
    the leakage--noise tradeoff as a function of tilt range.
  \item Compute $\Hlm(R)$ and verify $C$-matrix rank for the
    planned tilt range.
  \item Implement the constructive estimator (Steps 1--2) on
    synthetic data and verify exact recovery under ideal conditions.
  \item Implement and test the ALS/Gibbs refinement initialized from
    the constructive estimator.
  \item Run the prior sweep: compute $\sigma^\mathrm{post}_\ell$
    vs.\ $\sigma^\mathrm{prior}_\ell$ for the EIGSEP rotation grid
    and noise budget.
    Identify the $\ell$-modes where rotation data are informative
    beyond the HFSS simulation, and the minimum prior width at
    which the data dominate.
  \item Quantify the $M^{-1}\delta M$ transfer matrix for realistic
    LiDAR-derived horizon errors at the EIGSEP site.
  \item Determine the minimum tilt range for full $C$-matrix rank
    as a function of $\lmax$ for a flat and a non-flat horizon.
\end{enumerate}

\bibliography{refs}{}
\bibliographystyle{aasjournal}

\end{document}